% ==============================================================================
% capitulos/03-marco-normativo.tex - Marco Teórico sobre la normatividad y procesos de la DDP
% ==============================================================================

\chapter{Marco Teórico}
El presente capitulo establece los fundamentos normativos, operativos y tecnológicos que respaldan el desarrollo de la plataforma. A través de este análisis, se justifican los procesos institucionales con las soluciones de ingeniería que se proponen, validando que el desarrollo tecnológico además de resolver la problemática operativa, este abarca las bases funcionales y rigor técnico que un proyecto de esta índole demanda en el ámbito profesional.

\section{Contexto Institucional y Normativo}
En esta sección se definen las bases legales del órgano para el cual se desarrolla la plataforma.

\subsection{Defensoría de Derechos Politécnicos (DDP)}
La DDP se establece como un órgano dotado de autonomía técnica, cuya misión fundamental es recibir y gestionar quejas relativas a actos u omisiones que vulneren o invaliden los derechos individuales otorgados a los miembros pertenecientes a la comunidad del Instituto Politécnico Nacional. Su creación responde a la necesidad de contar con una instancia independiente que garantice la legalidad y respeto a los derechos humanos dentro de la institución [3].

De acuerdo con su marco normativo, la Defensoría tiene la facultad de proponer soluciones a través de recomendaciones y mediaciones, siempre bajo los principios de imparcialidad y confidencialidad [3].

\subsection{Manuales de Organización y Procedimientos}
La operatividad de la plataforma web se rige bajo la estructura y los procesos definidos en los manuales administrativos oficiales de la Defensoría:
\begin{itemize}
    \item \textbf{Manual de Organización de la DDP:} Este documento determina la estructura jerárquica y las atribuciones específicas de los servidores públicos designados a la defensoría. Es el pilar para la definición del Control de Acceso Basado a Roles (RBAC) del sistema, permitiendo diferenciar las capacidades de cada rol, que tiene acciones en el proceso de la queja. [4].
    \item \textbf{Manual de Procedimientos de la DDP:} Este documento establece la secuencia de los procesos y los requisitos formales de cada etapa de atención de una queja, desde su recepción hasta la resolución final. Este manual dicta las reglas de negocio que el software debe automatizar, incluyendo los criterios de validación de información y generación de documentos oficiales que integran el expediente digital [5].
\end{itemize}

\subsection{Protocolo para la Prevención, Detección, Atención y Sanción de la Violencia de Género}
Además de la normativa orgánica, el actuar de la Defensoría está delimitado por el "Protocolo para la Prevención, Detección, Atención y Sanción de la Violencia de Género en el Instituto Politécnico Nacional" [12], el cual establece las rutas institucionales de actuación ante manifestaciones de violencia o discriminación en la comunidad politécnica.

Su relevancia legal radica en que define los límites de competencia de la Defensoría al diferenciar conductas administrativas (sancionables internamente) de posibles delitos. Con base en esto, el protocolo obliga a la institución a dictar medidas de protección y orientar a la persona quejosa para interponer su denuncia ante el Ministerio Público o la Fiscalía cuando la gravedad de los hechos excede la jurisdicción puramente administrativa del IPN.

\section{Lógica de Negocio y Flujo de Procesos}
Esta sección define los conceptos operativos y las reglas de negocio que sigue el funcionamiento del sistema, basándose en las secuencias de actividades y las estructuras jerárquicas dictadas por los manuales institucionales vigentes [4][5] y el acuerdo de creación de la DDP [3]. Al modelar estos procesos, se garantiza que la plataforma este apegada a las normativas y salvaguarde la integridad de cada etapa del proceso, desde el ingreso de la solicitud hasta la conclusión del expediente.

\subsection{Flujo Operativo de la Queja}
El proceso de atención a una queja dentro de la Defensoría de los Derechos Politécnicos se rige por un flujo operativo estandarizado que garantiza la atención jurídica desde el primer contacto hasta el archivo del expediente. De acuerdo con el seguimiento de juntas con la DDP y en alineación con el Manual de Procedimientos, este ciclo se compone de las siguientes fases normativas:

\begin{itemize}
    \item \textbf{Recepción y validación inicial de la solicitud:} El proceso inicia cuando la Persona Quejosa presenta su queja ante la Defensoría. La Oficina de Primer Contacto recibe y analiza la documentación para revisar que esté completa. En esta etapa se valida la mayoría de edad; en caso de ser menor, se solicita la ficha de datos del tutor y contacto. Si la documentación cumple con los requisitos, se coloca el sello de recibido y se asigna un número de control inicial; en caso contrario, se solicita a la persona que complemente la queja o se emite un mensaje de rechazo.
    \item \textbf{Determinación del trámite:} Una vez conformada la documentación inicial, el expediente es turnado al Titular de la DDP. Esta figura de autoridad recibe, analiza la petición y determina la atención, así como el trámite correspondiente, derivando las instrucciones hacia la Subdefensoría para su seguimiento.
    \item \textbf{Análisis de competencia y Radicación:} Con la instrucción del Titular, las áreas jurídicas determinan si existe competencia institucional sobre los hechos narrados, evaluando la presencia de antecedentes y apoyándose en normativas como el Protocolo de Violencia de Género. De confirmarse la procedencia, se elabora el acuerdo de admisión y se asigna el número de expediente definitivo.
    \item \textbf{Investigación y solicitud de medidas:} Radicado el expediente, la Subdefensoría determina si las características del caso requieren elaborar una solicitud de medidas de protección para salvaguardar a la persona afectada. De forma paralela, se elaboran y envían oficios solicitando actos de investigación a los directores de las unidades correspondientes, verificando periódicamente que haya respuesta y enviando recordatorios cuando es necesario.
    \item \textbf{Resolución y conclusión:} Tras el análisis de los informes recibidos de las autoridades y la valoración de las pruebas, la Subdefensoría determina si existen los elementos suficientes para concluir el caso. Finalmente, se elabora el acuerdo de conclusión (pronunciamiento, recomendación o cese) y su respectiva notificación a la persona quejosa, turnando el caso al área secretarial para su archivo.
\end{itemize}

\subsection{Roles Administrativos y Operativos}
La estructura de los usuarios del sistema se fundamenta en la organización jerárquica y las facultades definidas en el Manual de Organización [4]. Para garantizar la integridad de la información y el cumplimiento de las responsabilidades legales, se definen los siguientes perfiles operativos:
\begin{itemize}
    \item \textbf{Titular de la Defensoría:} Responsable de dirigir, coordinar y supervisar el funcionamiento de la Defensoría, así como de emitir las recomendaciones y pronunciamientos finales derivados de los expedientes [4].
    \item \textbf{Cuerpo de Abogados Asesores (Subdefensoría):} Son los encargados de la atención directa de los casos, realizando el análisis jurídico, la integración de expedientes, la valoración de pruebas y la elaboración de proyectos de resolución [4].
    \item \textbf{Área Secretarial:} Área encargada de coordinar el apoyo administrativo y el seguimiento de acuerdos, siendo el enlace entre las diversas áreas de la Defensoría [4].
    \item \textbf{Personal de Apoyo (Primer Contacto):} Área responsable de la recepción inicial de las solicitudes, la captura de datos generales del quejoso y la gestión de la correspondencia oficial [4].
\end{itemize}

\subsection{Términos y Plazos dentro del Proceso}
En el ámbito administrativo, los términos procesales representan los lapsos temporales obligatorios en los que la autoridad debe realizar una actuación jurídica especifica [5]. El sistema debe integrar el monitoreo de estos plazos para asegurar que la gestión de la queja se mantenga dentro de los límites establecidos en la legalidad de la normativa institucional:
\begin{itemize}
    \item \textbf{Plazo para la Radicación:} Una vez recibida y analizada la solicitud, la Defensoría debe proceder a la radicación y dar turno del expediente en un periodo de 2 días habiles para evitar el rezago administrativo [5].
    \item \textbf{Término para la Solicitud de Informe:} Tras la admisión de la queja, se debe notificar a la autoridad responsable, otorgándole un plazo de (2 o 15) días para rendir su informe detallado sobre los hechos [5].
    \item \textbf{Periodo de Valoración:}  El abogado asesor cuenta con un tiempo definido para analizar las evidencias y los informes recibidos antes de proponer un proyecto de resolución [5].
    \item \textbf{Plazos de Seguimiento:} Una vez emitida una recomendación o acuerdo, el sistema debe vigilar los tiempos otorgados para el cumplimiento de las acciones remediales antes del archivo del caso, dando al quejoso 5 días hábiles para contestar lo que a derecho le convenga [5].
\end{itemize}

\subsection{Instrumentos de Captura y Gestión de Documentos}
La sistematización de las quejas requiere de la estandarización de los instrumentos de captura para garantizar que la información recolectada sea íntegra y suficiente para el análisis jurídico. Por lo que el sistema debe contemplar la digitalización y gestión de los siguientes documentos:
\begin{itemize}
    \item \textbf{Formato de Registro de Queja:} Este es el documento base que debe contener los datos de identificación del quejoso (nombre, unidad académica, número de boleta o empleado) y la narración de los hechos [5].
    \item \textbf{Catálogo de Derechos Vulnerados:} Herramienta que permite categorizar la queja según la normativa del IPN, facilitando la clasificación automática y la generación de estadísticas [5].
    \item \textbf{Acuse de Recibo y Folio de Seguimiento:} Folio generado automáticamente por el sistema que brinda certeza jurídica al usuario sobre el inicio de su trámite y los medios para consultar su estatus [5].
    \item \textbf{Expediente Digital de Evidencias:} Repositorio centralizado para el almacenamiento de archivos multimedia (imágenes, audios, videos) y documentos en formato PDF que sustentan la queja presentada [5].
    \item \textbf{Oficios de Notificación y Solicitud de Informe:} Formatos preestablecidos que el sistema debe emitir para requerir información a las autoridades señaladas como responsables, cumpliendo con las formalidades del Manual de Procedimientos [5].
\end{itemize}
